\section{Research Gap}
\label{sec:introduction_research_gap}

Research on virtual embodiment has shown that multisensory congruence, visuomotor synchrony, and avatar representation can influence how users experience virtual bodies and actions.
% TODO cite: embodiment / virtual body ownership / agency literature.

A separate body of work has explored virtual reality and interactive systems for sport, exercise, and rowing. These systems often aim to support training, motivation, feedback, or presence by augmenting stationary exercise with virtual environments, avatars, or simulated sport contexts.
% TODO cite: VR exercise / VR sport training review.
% TODO cite: rowing VR / rowing simulator sources.

Existing rowing-related VR and simulation work already contains many of the ingredients relevant to this thesis. Some systems use ergometer like vehicles as virtual representations. Others use 3D models of rowing machines or rowing machine inspired vehicles to represent the User. Some systems use inverse kinematics or avatar animation to represent the user's rowing movement. Commercial and consumer-facing applications also illustrate different representational choices, ranging from travelling through virtual environments on an indoor machine to more boat- or oar-centered rowing depictions.
% TODO cite: rowing VR systems / SPRINT / ETH rowing simulator / Holofit or other commercial examples if appropriate.
% TODO cite: attached evidence review or source map if converted into thesis note.

However, the specific comparison addressed in this thesis appears to be missing from the reviewed literature. Existing studies tend to compare VR against non-VR rowing, evaluate feedback or visualization techniques, manipulate avatar properties, adjust simulator parameters, or study motivational and social factors. They generally do not isolate the physical-to-virtual rowing mapping itself as the central experimental variable.

More specifically, the reviewed material does not appear to contain a controlled study that holds the physical rowing setup constant while comparing an ergometer-based mapping against a boat-based mapping as the main independent variable. Such a comparison would ask whether the avatar and rowing apparatus should be solved primarily in the coordinate frame of the real ergometer, or transformed into the coordinate frame of a virtual boat and oar system.
% TODO cite: evidence review / literature map showing no direct head-to-head comparison found.
% TODO cite: examples of studies that manipulate avatar, feedback, VR vs non-VR, simulator parameters.

This gap is important because the mapping choice is not only a technical implementation detail. It changes the relation between the user's body, the physical machine, and the virtual action they see. As a result, it may influence agency, body ownership, self-location, perceived realism, and preference. If the mapping is not isolated experimentally, it remains unclear whether a more sport-realistic rowing depiction supports the experience more strongly than a more sensorimotor-congruent ergometer depiction, or whether the opposite is the case.

The thesis addresses this gap by focusing on the mapping choice itself. It compares an ergometer-congruent representation and a boat-centric representation within the same general VR rowing context. The aim is not to evaluate every aspect of VR rowing, but to examine how this specific design choice shapes embodied and experiential responses.
